\documentclass[10pt]{article}
\usepackage[T1]{fontenc}
\usepackage[utf8]{inputenc}
\usepackage[french]{babel}
\usepackage{amsmath,amssymb,mathrsfs,mathpazo}
\usepackage{geometry}
\geometry{a4paper,margin=2cm}
\usepackage{graphicx,array,tabularx,tabularray,colortbl,multirow,multicol,booktabs,enumitem,pifont,ulem,fancybox,diagbox,cancel}
\usepackage{pstricks,pstricks-add,pst-plot,pst-tree,pst-node,pst-text,pst-func,pst-math,pst-bspline,pst-3dplot,pst-eucl,pst-solides3d}
\let\pscalculate\undefined
\usepackage{pst-bezier}
\usepackage{tikz,pgfplots,tkz-tab}
\usetikzlibrary{arrows,arrows.meta,patterns,calc,positioning,trees,decorations.markings,shapes.geometric,shapes.misc,shapes.arrows,intersections,decorations.pathmorphing,decorations.pathreplacing,backgrounds,shadows,babel}
\pgfplotsset{compat=1.18}
\usepgfplotslibrary{fillbetween}
\usepackage[np]{numprint}
\usepackage{listings,makecell,pas-tableur,twoopt}
\let\barre\undefined
\usepackage[pstricks]{tablvar}
\definecolor{darkgreen}{RGB}{0,100,0}
\providecommand{\red}{\color{red}}
\providecommand{\blue}{\color{blue}}
\providecommand{\green}{\color{green}}
\providecommand{\black}{\color{black}}
\providecommand{\white}{\color{white}}
\providecommand{\gray}{\color{gray}}
\providecommand{\cyan}{\color{cyan}}
\providecommand{\magenta}{\color{magenta}}
\providecommand{\yellow}{\color{yellow}}
\providecommand{\violet}{\color{violet}}
\providecommand{\orange}{\color{orange}}
\providecommand{\euro}{\texteuro}
\providecommand{\EUR}{\texteuro}
\providecommand{\ieme}{\textsuperscript{e}}
\providecommand{\degres}{\ensuremath{^\circ}}
\providecommand{\degree}{\ensuremath{^\circ}}
\providecommand{\vv}[1]{\overrightarrow{#1}}
\providecommand{\Coord}[1]{#1}
\providecommand{\decosix}{}
\providecommand{\starredbullet}{\ensuremath{\bullet}}
\pagestyle{empty}
\setlength{\parindent}{0pt}
\usepackage{adjustbox}\DeclareUnicodeCharacter{2236}{:}\DeclareUnicodeCharacter{2190}{\ensuremath{\leftarrow}}
\providecommand{\R}{\mathbb{R}}
\providecommand{\N}{\mathbb{N}}
\providecommand{\Z}{\mathbb{Z}}
\providecommand{\Q}{\mathbb{Q}}
\providecommand{\D}{\mathbb{D}}
\providecommand{\C}{\mathbb{C}}
\providecommand{\e}{{\mathrm{e}}}
\providecommand{\vect}[1]{\overrightarrow{#1}}
\providecommand{\vectt}[1]{\overrightarrow{\mathrm{#1}}}
\providecommand{\vvt}[1]{\overrightarrow{\mathrm{#1}}}
\providecommand{\np}[2][]{{#2~#1}}
\providecommand{\vv}[1]{{\overrightarrow{#1}}}
\renewcommand{\d}{{\mathrm{d}}}
\providecommand{\degres}{{^\circ}}
\providecommand{\degre}{{^\circ}}
\providecommand{\blue}{\color{blue}}
\providecommand{\red}{\color{red}}
\providecommand{\enskip}{\quad}
\providecommand{\llbracket}{\lbrack\!\lbrack}
\providecommand{\rrbracket}{\rbrack\!\rbrack}
\providecommand{\euro}{{\text{€}}}
\providecommand{\no}{{\text{n°}}}
\providecommand{\interval}[2]{{[#1;#2]}}
\providecommand{\intervFO}[2]{{[#1;#2[}}
\providecommand{\dotfill}{\ldots}

\providecommand{\ds}{\displaystyle}
\providecommand{\textding}[1]{#1}
\providecommand{\Oij}{\(\mathrm{O};\vec{\imath},\vec{\jmath}\)}
\providecommand{\Oijk}{\(\mathrm{O};\vec{\imath},\vec{\jmath},\vec{k}\)}
\providecommand{\Ouv}{\(\mathrm{O};\vec{u},\vec{v}\)}
\providecommand{\intervFF}[2]{\ensuremath{[#1;#2]}}
\renewcommand{\euro}{\ifmmode\text{\texteuro}\else\texteuro\fi}
\DeclareUnicodeCharacter{00B0}{\ifmmode^{\circ}\else\textdegree\fi}
\DeclareUnicodeCharacter{20AC}{\ifmmode\text{\texteuro}\else\texteuro\fi}
% Énoncés : Éducation nationale. Transcription et corrigés : APMEP, Denis Vergès et contributeurs. Découpage et classement : CoopMaths, Sébastien Lozano, Éric Elter et contributeurs. Adaptation de la présentation : Maths Explorer.
% CC BY-SA : https://coopmaths.fr/www/about/
% https://coopmaths.fr/alea/static/bac/2022/tex/bac_2022_05_sujet2_madagascar_1.tex
\begin{document}
\noindent\textit{APMEP / CoopMaths. CC BY-SA. Présentation adaptée par Maths Explorer.}\par\medskip
\medskip

\emph{Les résultats seront arrondis si besoin à $10^{- 4}$ près}

\medskip

Une étude statistique réalisée dans une entreprise fournit les informations suivantes : 

\setlength\parindent{1cm}
\begin{itemize}
\item[$\bullet~~$] 48\,\% des salariés sont des femmes. Parmi elles, 16,5\,\% exercent une
profession de cadre ;
\item[$\bullet~~$] 52\,\% des salariés sont des hommes. Parmi eux, 21,5\,\% exercent une
profession de cadre.
\end{itemize}
\setlength\parindent{0cm}

\medskip

On choisit une personne au hasard parmi les salariés. On considère les évènements suivants:

\setlength\parindent{1cm}
\begin{itemize}
\item[$\bullet~~$] $F$: \og la personne choisie est une femme \fg{} ;
\item[$\bullet~~$] $C$: \og la personne choisie exerce une profession de cadre \fg.
\end{itemize}
\setlength\parindent{0cm}

\medskip

\begin{enumerate}
\item Représenter la situation par un arbre pondéré.
\item Calculer la probabilité que la personne choisie soit une femme qui exerce une profession de cadre.
\item 
	\begin{enumerate}
		\item Démontrer que la probabilité que la personne choisie exerce une profession de cadre est égale à $0,191$.
		\item Les évènements $F$ et $C$ sont-ils indépendants ? Justifier.
	\end{enumerate}	
\item Calculer la probabilité de $F$ sachant $C$, notée $P_C(F)$. Interpréter le résultat dans le
contexte de l'exercice.
\item On choisit au hasard un échantillon de $15$ salariés. Le grand nombre de salariés
dans l'entreprise permet d'assimiler ce choix à un tirage avec remise. 

On note $X$ la variable aléatoire donnant le nombre de cadres au sein de
l'échantillon de 15 salariés.

On rappelle que la probabilité qu'un salarié choisi au hasard soit un cadre est
égale à $0,191$.
	\begin{enumerate}
		\item Justifier que $X$ suit une loi binomiale dont on précisera les paramètres. 
		\item Calculer la probabilité que l'échantillon contienne au plus 1 cadre.
		\item Déterminer l'espérance de la variable aléatoire $X$.
	\end{enumerate}	
\item Soit $n$ un entier naturel.

On considère dans cette question un échantillon de $n$ salariés.

Quelle doit être la valeur minimale de $n$ pour que la probabilité qu'il y ait au moins un cadre au sein de l'échantillon soit supérieure ou égale à $0,99$ ?
\end{enumerate}

\bigskip

\end{document}
